\documentclass[hyphens,aspectratio=169,dvipsnames]{beamer}
\usepackage{graphicx}
\usepackage{xcolor}
\usepackage{amssymb}
\usepackage{amsmath}
\usepackage{mathtools}
\usepackage{textcomp}
\usepackage{moresize}
\usepackage{framed}
\usepackage{minted}
\usepackage{relsize}
\usepackage{tikz}
\usepackage{comment}
\usetikzlibrary{shapes.geometric, arrows, positioning}
\usetheme{Berlin}
\usecolortheme[RGB={0,95,47}]{structure}
\beamertemplatenavigationsymbolsempty

\makeatletter
\AtBeginEnvironment{minted}{\dontdofcolorbox}
\def\dontdofcolorbox{\renewcommand\fcolorbox[4][]{##4}}
\makeatother

\newcommand{\textpf}[1]{\texttt{\color{black}\fcolorbox{lightgray}{lightgray}{#1}}}

\begin{document}
\section{Binary and Hex}
\begin{frame}{Learning objectives}
    \begin{itemize}
        \item Understanding numbering basics
        \item How to count in different bases
        \item Understanding how signed and unsigned numbers are encoded
        \item Understanding argv, program arguments and arrays of pointers
    \end{itemize}
\end{frame}

\begin{frame}{Binary and Hexadecimal}
    \begin{itemize}
        \item What is binary? What base is that? When do we use it?
        \pause \item What base do humans usually use?
        \pause \item What is hexadecimal (often referred to as just ``hex'')? Why do we use it?
    \end{itemize}
    % Whiteboard as necessary here.
\end{frame}

\section{Let's do some counting!}

\begin{frame}{Counting in binary}
    \begin{itemize}
        \item Count from $2^0$ to $2^4$ in binary on paper.
        \item Same thing in hex.
            \pause \item Now let's do some additions and subtractions!
            \pause \item How much is $ff + 1$? (in base 16)
            \pause \item How much is $200 -1$ (in base 16)
    \end{itemize}
\end{frame}

\begin{frame}{Integer literals in C}
    In C, the base of an integer literal is specified by special prefixes:
    \begin{itemize}
        \pause \item \textpf{0b} indicates base 2
        \pause \item \textpf{0} indicates base 8
        \pause \item \textpf{0x} indicates base 16
        \pause \item (no prefix indicates base 10)
    \end{itemize}
\end{frame}

\begin{frame}{Activity: Arithmetic on Paper}
    Please compute the following (assume C integer literal syntax).
    \begin{itemize}
        \item \textpf{0b1010 * 2} (give answer in base 2)
        \item \textpf{0xc001 * 16} (give answer in base 16)
        \item \textpf{0b0101 + 0b1011} (give answer in base 2)
        \item \textpf{0xda1e + 0xdafe} (give answer in base 16)
        \item \textpf{0xc001 + 010} (give answer in base 16)
    \end{itemize}
\end{frame}

\section{2's Complement}

\begin{frame}{Negative Binary Numbers}
    How should we represent negative numbers in binary?
\end{frame}

\begin{frame}{Take 1: Sign Bit}
    \begin{itemize}
        \pause \item Reserve the most significant bit as the sign bit.
        \pause \item To negate a number, flip its sign bit.
    \end{itemize}

    \begin{align*}
        \texttt{bin}(-25) &= \texttt{set\_msb}(\texttt{bin}(25)) \\
                          &= \texttt{set\_msb}(\texttt{0b}00011001) \\
                          &= \texttt{0b}10011001
    \end{align*}
\end{frame}

\begin{frame}{Sign Bit Pros and Cons}
    Pros
    \begin{itemize}
        \pause \item It's easy for people to understand.
    \end{itemize}

    \pause

    Cons
    \begin{itemize}
        \pause \item +0 and -0 are distinct.
        \pause \item Calculations are not consistent with unsigned numbers.
    \end{itemize}

    \pause
    \begin{center}Whiteboard!\end{center}
\end{frame}

\begin{frame}{Take 2: 2's Complement}
    \begin{itemize}
        \pause \item To negate a number, flip all of its bits, then add 1.
        \pause \item This is how nearly all modern computers represent negative numbers.
    \end{itemize}

    \pause \begin{align*}
        \texttt{bin}(-25) &= \texttt{flip\_bits}(\texttt{bin}(25)) + \texttt{0b}1 \\
                          &= \texttt{flip\_bits}(\texttt{0b}00011001) + \texttt{0b}1 \\
                          &= \texttt{0b}11100110 + \texttt{0b}1 \\
                          &= \texttt{0b}11100111
    \end{align*}
\end{frame}

\begin{frame}{2's Complement Pros and Cons}
    Pros
    \begin{itemize}
        \pause \item There's only one zero.
        \pause \item Calculations are consistent with unsigned numbers!
    \end{itemize}

    \pause
    \begin{center}Whiteboard!\end{center} % -18 -38
    \pause

    Cons
    \begin{itemize}
        \pause \item It's more difficult to understand.
    \end{itemize}
\end{frame}


\section{\textpf{argc} and \textpf{argv}}

\begin{frame}[fragile]{\textpf{argc} and \textpf{argv}}
    So far, we've used \textpf{main} functions with the following declaration:
    \begin{minted}{c}
int main(void);
    \end{minted}
\end{frame}

\begin{frame}[fragile]{\textpf{argc} and \textpf{argv}}
    It's also typical to use this declaration:
    \begin{minted}{c}
int main(int argc, char **argv);
    \end{minted}
\end{frame}

\begin{frame}[fragile]{\textpf{argc} and \textpf{argv}}
    \begin{itemize}
        \pause \item \textpf{argv} is an array of strings containing the \textit{command line arguments} passed to this program.
        \pause \item \textpf{argc} is its length.
    \end{itemize}
\end{frame}

\begin{frame}[fragile]{\textpf{argv} is not \textpf{stdin}}
    \begin{itemize}
        \pause \item \textpf{argv} is the array of argument strings passed to the program on the command line \textit{before the program runs}.
        \pause \item \textpf{stdin} is a file descriptor that represents the input to the program \textit{at runtime}.
        \pause \item Example in terminal!
    \end{itemize}
\end{frame}


\begin{frame}[fragile]{Activity: Implement \textpf{echo}}
    \begin{itemize}
        \pause \item The \textpf{echo} program writes its arguments to \textpf{stdout}, separated by spaces.
        \pause \item e.g., \textpf{echo a b c} should write ``\textpf{a b c}" to \textpf{stdout}.
        \pause \item What do you notice about \textpf{argv[0]}?
    \end{itemize}
\end{frame}

\end{document}
